\chapter{Optimization}
\label{ch:optimization}

\section{Why optimize first}

Before asking how uncertain the parameters are, it is worth asking whether the model can represent
the data at all. Optimization answers that: minimize a loss $f(\theta)$ that measures the misfit,
for example $\chi^2(\theta)$, and check that its minimum is statistically acceptable. If it is not,
no amount of posterior sampling will help --- the problem lies with the model. If it is, the
minimizer has also located the region where the posterior lives, which is exactly where MCMC should
start its walkers (\cref{ch:mcmc}).

The problem solved is a bounded, optionally constrained minimization,
\begin{equation}
  \min_{\theta}\; f(\theta)
  \quad\text{subject to}\quad
  \theta \in [\ell_1, u_1] \times \dots \times [\ell_d, u_d],
  \qquad c(\theta) \le 0,
  \label{eq:minimize}
\end{equation}
where the box is \code{param_bounds} and $c$ is the optional \code{constraint_fun}: a point is
feasible when $c(\theta) \le 0$. The function $f$ is a black box: no gradients, and every value
costs an expensive evaluation. What can be parallelized is the number of points evaluated per
iteration, \code{num_workers}.

\section{The algorithms}

Three families are available. The first two are population methods from \texttt{nevergrad}; the
third builds a statistical model of $f$.

\subsection{Differential Evolution (default)}

Differential Evolution (DE) keeps a population of candidate points and improves it generation by
generation. For each member $\theta_i$ it builds a mutant from three other members,
\begin{equation}
  v_i = \theta_a + F\,(\theta_b - \theta_c),
\end{equation}
mixes coordinates of the mutant and of $\theta_i$ (crossover) into a trial point, and replaces
$\theta_i$ by the trial point if it has a lower loss. The difference vectors adapt the step size
to the population's own spread, which makes DE robust on rugged and badly scaled losses. The
default in slurmcmc is \texttt{nevergrad}'s DE with two-point crossover and a population equal to
\code{num_workers}, so one generation is exactly one iteration of parallel evaluations.

\subsection{Particle Swarm Optimization}

Particle Swarm Optimization (PSO) moves each particle with a velocity that is pulled towards the
best point the particle has seen, $p_i$, and the best point the swarm has seen, $g$:
\begin{equation}
  u_i \leftarrow \omega\, u_i + c_1 r_1 (p_i - \theta_i) + c_2 r_2 (g - \theta_i),
  \qquad
  \theta_i \leftarrow \theta_i + u_i,
\end{equation}
where $r_1$ and $r_2$ are random numbers in $[0,1]$. Select it with
\code{optimizer_class=ng.optimizers.ConfPSO(popsize=num_workers)}; any other \texttt{nevergrad}
optimizer can be passed the same way.

\subsection{Bayesian optimization}

Bayesian optimization (BO), with \code{optimizer_package='botorch'}, fits a Gaussian-process model
to all evaluations so far and uses it to decide where to evaluate next. The model gives, at every
point, a predicted loss and an uncertainty. A candidate is chosen to maximize the expected
improvement over the best loss found so far, $f^\ast$,
\begin{equation}
  \mathrm{EI}(\theta) \;=\; \mathbb{E}\!\left[\max\!\big(f^\ast - f(\theta),\, 0\big)\right],
\end{equation}
which balances exploiting regions of low predicted loss against exploring regions of high
uncertainty. slurmcmc uses \texttt{botorch}'s batch version, \code{qLogExpectedImprovement}, to
propose \code{num_workers} points per iteration, with inputs normalized and outputs standardized;
the first iteration, with no data yet, uses a quasi-random Sobol design.

BO spends far fewer evaluations when each one is very expensive, but the model has a cost of its
own. Fitting a Gaussian process to $n$ points scales as $\mathcal{O}(n^3)$, and it is refitted every
iteration, so a long run slows down. Two options in \code{botorch_kwargs} address this:

\begin{center}
\small
\begin{tabularx}{\linewidth}{@{}l l L@{}}
  \toprule
  \textbf{key} & \textbf{default} & \textbf{meaning} \\
  \midrule
  \code{num_best_points} & \code{None} & train the model only on this many lowest-loss points, keeping each fit bounded in size \\
  \code{sequential}      & \code{True} & choose the batch greedily, one point at a time, rather than optimizing all $q\times d$ coordinates jointly; much faster at equivalent quality \\
  \code{num_restarts}    & \code{10}   & restarts of the acquisition-function optimization \\
  \code{raw_samples}     & \code{100}  & random samples seeding those restarts \\
  \code{options}         & \code{None} & extra options for \texttt{botorch}'s \code{optimize_acqf} \\
  \bottomrule
\end{tabularx}
\end{center}

\section{How a batch of candidates is chosen}

Each iteration, \code{slurm_minimize} asks the optimizer for candidates until it holds
\code{num_workers} of them that are
\begin{itemize}
  \item \textbf{new} --- not evaluated in any earlier iteration, nor already in this batch;
  \item \textbf{feasible} --- satisfying \code{constraint_fun} $\le 0$, which is evaluated cheaply
        before any expensive job is spent on the point.
\end{itemize}
Only then is the batch submitted as one job array. If \code{num_asks_max} requests pass without
completing the batch --- typically because the feasible region is tiny --- the run stops with an
error rather than looping forever.

Starting points can be supplied as \code{init_points}; they are evaluated in the first iteration
and must satisfy the constraint. An iteration in which \emph{every} evaluation failed stops the
run, since the optimizer would have nothing to learn from.

\needspace{14\baselineskip}
\section{Usage}

\begin{lstlisting}
import numpy as np
from slurmcmc.optimization import slurm_minimize

def loss_fun(x):
    return (1 - x[0])**2 + 100 * (x[1] - x[0]**2)**2

def constraint_fun(x):                      # feasible when <= 0
    return (x[0] + 1)**2 + (x[1] + 1)**2 - 16

status = slurm_minimize(loss_fun=loss_fun, constraint_fun=constraint_fun,
                        param_bounds=[[-5, 5], [-5, 5]],
                        num_workers=10, num_iters=30,
                        cluster='slurm', work_dir='minimize', save_restart=True)
print(status['x_min'], status['loss_min'])
\end{lstlisting}

The arguments specific to minimization are listed below; those shared by every function are in
\cref{app:common}.

\begin{argtable}
  \code{loss_fun} & required & the loss, \code{f(x)} or \code{f(x, extra_arg)}; a callable or a deferred-import dictionary \\
  \code{param_bounds} & required & \code{[[lo, hi], ...]}, one pair per parameter \\
  \code{num_workers} & required & points evaluated in parallel per iteration \\
  \code{num_iters} & required & iterations to run (further iterations, when resuming) \\
  \code{optimizer_package} & \code{'nevergrad'} & \code{'nevergrad'} or \code{'botorch'} \\
  \code{optimizer_class} & \code{None} & a \texttt{nevergrad} optimizer; \code{None} is Differential Evolution \\
  \code{botorch_kwargs} & \code{None} & options of Bayesian optimization, see the table above \\
  \code{init_points} & \code{None} & points to evaluate in the first iteration; must satisfy the constraint \\
  \code{constraint_fun} & \code{None} & \code{c(x)}, feasible when $\le 0$; checked before a point is evaluated \\
  \code{num_asks_max} & \code{1000} & optimizer requests allowed per iteration while filling the batch \\
  \code{job_fail_value} & NaN & result recorded for a failed evaluation \\
\end{argtable}

The returned status dictionary contains the best point \code{x_min} and its loss \code{loss_min};
the best loss up to each iteration, \code{loss_min_all_iter}, and within each iteration,
\code{loss_min_per_iter}; where the minimum was found, \code{loc_point_min_all_iter}; the
\code{slurm_pool} with the complete evaluation history; the \code{optimizer} itself; and counters
of loss, constraint and optimizer calls.

\section{Example: a constrained 2D Rosenbrock}
\label{sec:opt-example}

The loss is the two-dimensional Rosenbrock function inside a circle,
\begin{equation}
  \min_{x,y}\; (1-x)^2 + 100\,(y-x^2)^2
  \quad\text{subject to}\quad (x+1)^2 + (y+1)^2 \le 16,
\end{equation}
a narrow curved valley whose minimum at $(1,1)$ lies inside the constraint. Differential Evolution
with default settings runs for 30 iterations of 10 parallel evaluations.

\begin{figure}[H]
  \centering
  \includegraphics[width=0.78\linewidth]{example_optimization_loss_progress.png}
  \caption{Loss per iteration: all sampled points (blue), the best point of each iteration
           (green), and the best point found so far (red).}
  \label{fig:opt-loss}
\end{figure}

\begin{figure}[H]
  \centering
  \includegraphics[width=0.72\linewidth]{example_optimization_2d_visualization.png}
  \caption{The loss (logarithm of its absolute value), the circular constraint (white), and the
           sampled points, from dark (early) to bright (late), closing in on the minimum (star).}
  \label{fig:opt-2d}
\end{figure}

The number of iterations here is arbitrary. In a real calibration the loss measures the fit to data,
and the goal is to reach a statistically acceptable value before moving on to MCMC.

\section{Comparing the algorithms}
\label{sec:opt-comparison}

The three algorithms are compared on the $d$-dimensional Rosenbrock function,
\begin{equation}
  f(\theta) = \sum_{i=1}^{d-1} \Big[ 100\,(\theta_{i+1} - \theta_i^2)^2 + (1-\theta_i)^2 \Big],
  \label{eq:rosenbrock}
\end{equation}
whose minimum $f = 0$ lies at $\theta = (1, \dots, 1)$, over $[-10, 10]^d$ without constraints.
Every run uses 50 parallel evaluations for 50 iterations. DE and PSO use \texttt{nevergrad}'s
defaults; BO uses \code{num_restarts=10}, \code{raw_samples=100}, and no trimming of its training
set. For a fair comparison all three start from the same 50 random points, and the whole comparison
is repeated for three different sets of starting points (line styles). The legend gives the run
time of each full optimization.

\begin{figure}[H]
  \centering
  \includegraphics[width=0.74\linewidth]{example_optimization_algorithms_comparison_2d.png}
  \caption{Best loss per iteration for $d=2$.}
\end{figure}

\begin{figure}[H]
  \centering
  \includegraphics[width=0.74\linewidth]{example_optimization_algorithms_comparison_5d.png}
  \caption{Best loss per iteration for $d=5$.}
\end{figure}

\begin{figure}[H]
  \centering
  \includegraphics[width=0.74\linewidth]{example_optimization_algorithms_comparison_10d.png}
  \caption{Best loss per iteration for $d=10$.}
\end{figure}

\begin{keypoint}[What the comparison shows]
On these losses PSO outperforms both DE and BO. BO costs disproportionately more computation per
iteration --- its Gaussian process is refitted on a growing data set --- while rarely improving on
the alternatives, briefly at $d=10$ being the exception. BO earns its cost only when each evaluation
is so expensive that saving evaluations dominates everything else.
\end{keypoint}
